\documentclass[10pt,letterpaper]{article}

\usepackage[T1]{fontenc}
\usepackage[utf8]{inputenc}
\usepackage{charter}
\usepackage[
  letterpaper,
  left=0.68in,
  right=0.68in,
  top=0.52in,
  bottom=0.55in,
  includefoot
]{geometry}
\usepackage{xcolor}
\usepackage{tabularx}
\usepackage{array}
\usepackage{fancyhdr}
\usepackage{hyperref}

\definecolor{accent}{HTML}{843542}
\definecolor{bodytext}{HTML}{292525}
\definecolor{muted}{HTML}{665F5F}
\definecolor{rulecolor}{HTML}{CBBDB5}

\hypersetup{
  colorlinks=true,
  urlcolor=accent,
  linkcolor=accent,
  pdftitle={Curriculum Vitae - Johannes Hosle},
  pdfauthor={Johannes Hosle}
}

\color{bodytext}
\setlength{\parindent}{0pt}
\setlength{\parskip}{0pt}
\setlength{\tabcolsep}{0pt}
\setlength{\emergencystretch}{2em}
\linespread{0.96}

\pagestyle{fancy}
\fancyhf{}
\fancyfoot[L]{\color{muted}\footnotesize Johannes Hosle}
\fancyfoot[R]{\color{muted}\footnotesize\thepage}
\renewcommand{\headrulewidth}{0pt}
\renewcommand{\footrulewidth}{0pt}
\setlength{\footskip}{0.31in}

\makeatletter
\newcommand{\footerline}{%
  \begingroup
  \color{rulecolor}\rule{\textwidth}{0.35pt}\par
  \vspace{5pt}%
  \endgroup
}
\fancyfoot[C]{\raisebox{11pt}[0pt][0pt]{\makebox[\textwidth][l]{\footerline}}}
\makeatother

\newcommand{\cvsection}[1]{%
  \vspace{4pt}%
  {\color{accent}\bfseries\fontsize{10.5}{11.5}\selectfont\MakeUppercase{#1}}\par
  \vspace{1pt}%
  {\color{rulecolor}\hrule height 0.35pt}%
  \vspace{2.5pt}%
}

\newcommand{\datedheading}[2]{%
  \par\noindent
  \begin{minipage}[t]{0.81\textwidth}\vspace{0pt}\bfseries #1\end{minipage}%
  \hfill
  \begin{minipage}[t]{0.17\textwidth}\vspace{0pt}\raggedleft\color{muted}#2\end{minipage}%
  \par
}

\newcommand{\entry}[3]{%
  \datedheading{#1}{#2}\par
  \if\relax\detokenize{#3}\relax\else#3\par\fi
  \vspace{1pt}%
}

\newcommand{\publication}[1]{%
  \begingroup
  \hangindent=0.8em
  \hangafter=1
  #1\par
  \endgroup
  \vspace{2pt}%
}

\begin{document}
\thispagestyle{fancy}

{\fontsize{26}{29}\selectfont Johannes Hosle}\par
\vspace{1pt}
{\color{accent}\href{mailto:jhosle@mit.edu}{jhosle@mit.edu}}\par

\cvsection{Employment}
\entry{Carnegie Mellon University in Qatar}{2026--present}{%
  Postdoctoral Fellow in Mathematics\par
  Mentors: Layan El Hajj and Henrik Shahgholian
}

\cvsection{Education}
\entry{Massachusetts Institute of Technology}{2020--2026}{%
  Ph.D. in Mathematics\par
  Advisor: David Jerison
}
\entry{University of California, Los Angeles}{2016--2020}{%
  B.Sc. in Mathematics with Departmental Honors, \textit{summa cum laude}\par
  M.A. in Mathematics
}

\cvsection{Publications and Preprints}
\publication{J. Hosle, \href{https://arxiv.org/abs/2609.05107}{Monotonicity formulas in positively curved settings with applications to two-phase free boundary problems}, submitted.}
\publication{J. Hosle, \href{https://arxiv.org/abs/2607.10881}{The isomorphic section-projection problem for convex bodies}, submitted.}
\publication{J. Hosle and P. Ivanisvili, \href{https://arxiv.org/abs/2606.25350}{Geometric Block Exponents and a Uniform Mixed-Alphabet Sumset Inequality}, submitted.}
\publication{J. Hosle, \href{https://arxiv.org/abs/2606.07946}{Stolarsky-Type Inequalities in a Max-Convolution Problem}, submitted.}
\publication{J. Hosle, A. V. Kolesnikov, and G. V. Livshyts, On the $L_p$-Brunn--Minkowski and dimensional Brunn--Minkowski conjectures for log-concave measures, \textit{J. Geom. Anal.} (2021).}
\publication{J. Hosle, On extensions of the Loomis--Whitney inequality and Ball's inequality for concave, homogeneous measures, \textit{Adv. Appl. Math.} (2020).}
\publication{J. Hosle, On the Comparison of Measures of Convex Bodies via Projections and Sections, \textit{Int. Math. Res. Not. IMRN} (2021).}

\cvsection{Honors and Awards}
\entry{Sigma Xi}{2026}{}
\entry{Dean of Science Fellowship, MIT}{2020--2023}{}
\entry{Phi Beta Kappa}{2020}{}
\entry{Sherwood Prize, UCLA}{2020}{}
\entry{Dean's Honors List}{2016--2020}{}
\entry{Putnam Competition, Top 250}{2016}{}
\entry{National Merit Scholarship Semifinalist}{2015}{}

\cvsection{Teaching and Grading}
\entry{18.095 Mathematics Lecture Series, MIT}{Jan. 2026}{Teaching Assistant (J. Dunkel)}
\entry{18.065 Matrix Methods in Data Analysis, Signal Processing, and Machine Learning, MIT}{Spring 2025}{Teaching Assistant (Z. Chen)}
\entry{18.100B Real Analysis, MIT}{Fall 2024}{Course Assistant (L. Guth)}
\entry{18.100P Real Analysis, MIT}{Spring 2024}{Recitation Leader (K. Naff)}
\entry{18.01 Calculus, MIT}{Fall 2023}{Recitation Leader (L. Guth)}

\newpage
\thispagestyle{fancy}

\cvsection{Mentorship Experience}
\entry{Academic Mentor, Mathroots @ MIT}{July 2026}{Two-week mathematical talent accelerator for high school students hosted by MIT PRIMES.}
\entry{MIT Directed Reading Program}{Jan. 2025}{Mentored undergraduates Ryan Dong and Caleb Pascale in convex geometry, including the Busemann--Petty problem and Bourgain's bound for the isotropic constant.}

\cvsection{Presentations}
\entry{Analysis Seminar, Columbia University}{Sept. 2026}{Upcoming}
\entry{Analysis Seminar, Brown University}{Sept. 2026}{Upcoming}
\entry{Synergies between Geometry, Probability, and Computation in High Dimensions}{Aug. 2026}{ICERM, Brown University, Providence, Rhode Island}
\entry{Workshop in PDE and Applied Math in the Northeast Region}{Mar. 2026}{University of Connecticut, Storrs, Connecticut}
\entry{Pure Math Graduate Student Seminar, MIT}{Feb. 2025}{}
\entry{BIRS Workshop: Geometric Tomography}{Feb. 2020}{BIRS, Banff, Canada}
\entry{Workshop in Convexity and Geometric Aspects of Harmonic Analysis}{Dec. 2019}{Georgia Institute of Technology, Atlanta, Georgia}

\cvsection{Research Programs and Summer Schools}
\entry{Interactions between Convex Geometry and Spectral Analysis}{July 2025}{Universit\'e de Montr\'eal, Montr\'eal, Canada}
\entry{Particle Interactive Systems: Analysis and Computational Methods}{June 2024}{SLMath, Berkeley, California}
\entry{Georgia Tech REU in Mathematics}{Summer 2018}{Supervisors: Michael Lacey and Galyna Livshyts}

\cvsection{Programming}
\textbf{Python:} Proficient; background in statistical learning theory; experience with NumPy and scikit-learn.\par\vspace{4pt}
\textbf{Julia:} Proficient; used for gradient flows in free boundary problems, supersolution construction, and numerical estimates for weighted Wirtinger inequalities.\par\vspace{4pt}
\textbf{Wolfram Language:} Intermediate; familiar with machine learning methods.\par\vspace{4pt}
\textbf{LaTeX:} Proficient.

\cvsection{Languages and Citizenship}
\textbf{Languages:} English (native), Russian (elementary).\par\vspace{4pt}
\textbf{Citizenships:} United States and Germany.

\end{document}
